\section{Proof Overview}
\label{sec:proof-idea}

% \rnote{The section names are odd; also too many sections. Its helpful to state what a section is going to do at the beginning of the section. For example, 9.1. }

\begin{theorem}\label{thm:lower-bound}
The Grothendieck constant satisfies
\[
        K_G\ge\frac{6\pi}{11} = 1.7136\ldots.
\]
\end{theorem}

A sketch of the proof is as follows. First, recall the inverse-series step in Krivine's rounding argument. If
\[
        H^{-1}(z)=\sum_{j\ge0}a_{2j+1}z^{2j+1},
\]
define the absolute-coefficient majorant $M_H(s):=\sum_{j\ge0}|a_{2j+1}|s^{2j+1}$. The standard tensor preprocessing gives the implication
\begin{equation}
        M_H(\gamma)\le1
        \quad\Longrightarrow\quad
        K_G\le\frac{\pi}{2\gamma}.
        \label{eq:inverse-majorant-bound}
\end{equation}
We refer to~\eqref{eq:inverse-majorant-bound} as the \emph{inverse-majorant bound}.

The conceptual starting point is the theorem of Naor and Regev that Krivine schemes are optimal \cite{naor2014krivine}. By reworking the construction in their proof, we show that the inverse-majorant bound is itself asymptotically optimal after allowing mixtures and increasing dimension: there are limiting mixed correlation functions with admissible parameters converging to $\pi/(2K_G)$. This transfer is established in Section~\ref{sec:majorant_optimal}.

It therefore remains to prove a universal upper bound on such admissible parameters. For arbitrary odd measurable sign functions $f,g:\R^d\to\{\pm1\}$, write
\[
        H_{f,g}(t)=b_1t+b_3t^3+b_5t^5+\cdots.
\]
The technical core of the argument is the affine coefficient inequality
\[
        b_3\ge 2b_1-\frac{11}{6}.
\]
Section~\ref{sec:preliminaries} sets up the Gaussian--Hermite formulation, and Section~\ref{sec:affine-strip} proves this inequality by decomposing $f$ and $g$ into their agreement and disagreement parts and reducing the required estimates to one-dimensional inequalities. Its linearity in $(b_1,b_3)$ is essential, because it makes the inequality stable under mixtures and coefficientwise limits.

Section~\ref{sec:strip_majorant} then shows that the affine coefficient inequality controls the inverse-majorant bound: every admissible value must satisfy $\gamma\le11/12$. Applying this barrier to the asymptotically optimal sequence constructed in Section~\ref{sec:majorant_optimal} gives
\[
        \frac{\pi}{2K_G}\le\frac{11}{12},
\]
which is exactly Theorem~\ref{thm:lower-bound}.

\section{Preliminaries on Gaussian Hermite Expansions}
\label{sec:preliminaries}

Let
\[
        X\sim N(0,I_d)
\]
be a standard Gaussian vector in $\R^d$. We define the \textit{inner product} as the Gaussian expectation:
\[
        \ip{F}{G}=\E[F(X)G(X)].
\]
We use the orthonormal 1D Hermites
\[
        \HermBasis_0=1,\qquad
        \HermBasis_1(s)=s,\qquad
        \HermBasis_2(s)=\frac{s^2-1}{\sqrt2},\qquad
        \HermBasis_3(s)=\frac{s^3-3s}{\sqrt6}.
\]

In $d$ dimensions the Hermite basis is indexed by multi-indices. If
\[
        \alpha=(\alpha_1,\ldots,\alpha_d),
        \qquad |\alpha|=\alpha_1+\cdots+\alpha_d,
\]
then
\[
        \HermBasis_\alpha(X)=\prod_{i=1}^d \HermBasis_{\alpha_i}(X_i).
\]
Any square integrable function has a Hermite expansion, written as
\[
        F=P_0F+P_1F+P_2F+P_3F+\ldots,
\]
where $P_iF$ is the orthogonal projection of $F$ onto the span of all $\HermBasis_\alpha$ with total degree $|\alpha|=i$.

In particular, the degree-one Hermites span
\[
        \cH_1=\operatorname{span}\{X_1,\ldots,X_d\}.
\]
Here,
\[
        P_1F=\ip{c_F}{X},
        \qquad
        c_F=\E[F(X)X]\in\R^d.
\]
Its norm is
\[
        \norm{P_1F}^2=\norm{c_F}_{\ell_2}^2.
\]

Similarly, $P_3F$ is the projection onto the full total degree-three Hermite span. We will not list all basis elements. The important point is that it includes pure cubic directions such as $\HermBasis_3(X_i)$ and mixed directions such as $\HermBasis_2(X_i)\HermBasis_1(X_j)$ and $\HermBasis_1(X_i)\HermBasis_1(X_j)\HermBasis_1(X_\ell)$.

For signs $f,g:\R^d\to\{\pm1\}$, if
\[
        H_{f,g}(t)=\frac\pi2\E[f(X)g(X'_t)]
\]
with $(X,X'_t)$ a pair of correlated standard Gaussian vectors, then Mehler's identity gives
\[
        H_{f,g}(t)=\frac\pi2\sum_{m\ge0}\ip{P_mf}{P_mg}t^m.
\]
Here
\[
        b_1=[t]H_{f,g}=\frac\pi2\ip{P_1f}{P_1g},
        \qquad
        b_3=[t^3]H_{f,g}=\frac\pi2\ip{P_3f}{P_3g}.
\]

For the rest of the paper, we will use the shorthand
\[
        \nu:=\sqrt{\frac{2}{\pi}}=\E|Z|,
        \qquad Z\sim N(0,1).
\]

We want to prove
\begin{equation}
        b_3\ge 2b_1-\frac{11}{6}.
        \label{eq:affine-strip}
\end{equation}

For example, the hyperplane partitions lie on this line. For $f=g=\sgn(X_1)$,
\[
        \norm{P_1f}=\nu,
        \qquad
        \norm{P_3f}^2=\frac{\nu^2}{6}.
\]
Thus
\[
        (b_1,b_3)=\left(1,\frac16\right),
\]
and this point lies exactly on the line since $1/6=2-11/6$.

Instead of proving the strip directly, it suffices to prove
\begin{equation}
        B_1(f,g)
        :=\norm{P_1f}\norm{P_1g}
          +\ip{P_1f}{P_1g}
          -\ip{P_3f}{P_3g}
        \le \frac{11}{6}\nu^2.
        \label{eq:b1-bound}
\end{equation}
Indeed, by Cauchy--Schwarz,
\[
        \norm{P_1f}\norm{P_1g}\ge \ip{P_1f}{P_1g}.
\]
So~\eqref{eq:b1-bound} implies
\[
        2\ip{P_1f}{P_1g}-\ip{P_3f}{P_3g}\le \frac{11}{6}\nu^2.
\]
Multiplying by $\pi/2$ and using $\nu^2=2/\pi$ gives~\eqref{eq:affine-strip}.

\section{Proof of the Affine Coefficient Inequality}
\label{sec:affine-strip}

\subsection{The agreement--disagreement substitution}
\label{sec:agreement-disagreement}

We define a change of variables based on where $f$ and $g$ agree or disagree. Set the functions $h,k$ to be
\[
        h=\frac{f+g}{2},\qquad k=\frac{f-g}{2}.
\]
Then
\[
        f=h+k,
        \qquad
        g=h-k.
\]
Since $f,g$ are signs,
\[
        h,k\in\{-1,0,1\},
        \qquad
        hk=0,
        \qquad
        h^2+k^2=1.
\]
Think of $h$ as the agreement part and $k$ as the disagreement part between $f$ and $g$. This change of variables considerably simplifies the proof.

Let
\[
        u=P_1h,
        \qquad
        v=P_1k.
\]
Then
\[
        P_1f=u+v,
        \qquad
        P_1g=u-v.
\]
Hence
\[
        \ip{P_1f}{P_1g}=\norm{u}^2-\norm{v}^2.
\]
Also
\[
        \norm{P_1f}\norm{P_1g}
        =\norm{u+v}\norm{u-v}
        \le \norm{u}^2+\norm{v}^2,
\]
using $ab\le(a^2+b^2)/2$ after expanding $\norm{u\pm v}^2$.

For degree three,
\[
        P_3f=P_3h+P_3k,
        \qquad
        P_3g=P_3h-P_3k,
\]
so
\[
        \ip{P_3f}{P_3g}=\norm{P_3h}^2-\norm{P_3k}^2.
\]
Putting these together,
\begin{align*}
        B_1(f,g)
        &\le
        \bigl(\norm{u}^2+\norm{v}^2\bigr)
        +\bigl(\norm{u}^2-\norm{v}^2\bigr)
        -\bigl(\norm{P_3h}^2-\norm{P_3k}^2\bigr)\\
        &=2\norm{P_1h}^2-\norm{P_3h}^2+\norm{P_3k}^2.
\end{align*}

So it is enough to prove
\begin{equation}
\boxed{\displaystyle
        D_1(h,k):=2\norm{P_1h}^2-\norm{P_3h}^2+\norm{P_3k}^2
        \le \frac{11}{6}\nu^2
}
\label{eq:hk-reduction}
\end{equation}

In order to prove this, we have to handle each of the three expressions $\norm{P_1h}^2$, $\norm{P_3h}^2$, $\norm{P_3k}^2$ carefully. In particular, we need to prove upper bounds on $\norm{P_1h}^2$, $\norm{P_3k}^2$ and a lower bound on $\norm{P_3h}^2$.

\begin{remark}
    Note that in the inequality before~\eqref{eq:hk-reduction} the term $\norm{v}^2=\norm{P_1k}^2$ conveniently cancels. This is a result of the choice of the strip $b_3\ge 2b_1-11/6$. Stronger lower bounds can be proved with the same method, but this is more complicated because one also has to bound $\norm{P_1k}^2$.
\end{remark}

\subsection{Reduction to one Gaussian coordinate}
\label{sec:rotation}

The linear part of $h$ has the form
\[
        P_1h(X)=\ip{c_h}{X}.
\]
The coefficient vector is $c_h\in\R^d$, and
\[
        \norm{P_1h}=\norm{c_h}_{\ell_2}.
\]
If $c_h\neq0$, rotate coordinates so the first coordinate points in the $c_h$ direction. After this rotation, and after renaming the coordinates, we write
\[
        X=(S,T),
        \qquad S\sim N(0,1),\quad T\sim N(0,I_{d-1}),
\]
with $S$ independent of $T$, and
\[
        P_1h(S,T)=\norm{P_1h}\,S.
\]
If $P_1h=0$, choose any direction as $S$.

Define
\[
        x:=\frac{\norm{P_1h}}{\nu}\in[0,1].
\]
The bound $x\le1$ follows from
\[
        \norm{P_1h}=\sup_{\norm{e}=1}\E[h(X)\ip{e}{X}]
        \le \E|Z|=\nu.
\]
After rotation,
\[
        \ip{h}{S}=\norm{P_1h}=\nu x.
\]

\subsection{The agreement terms}
\label{sec:agreement}

The agreement part of $D_1$ is $2\norm{P_1h}^2-\norm{P_3h}^2$. The helpful negative term is
\[
        -\norm{P_3h}^2.
\]
To use it, we only need a lower bound on $\norm{P_3h}^2$.

In any dimension,
\[
        \norm{P_3h}^2
        \ge \ip{h}{\HermBasis_3(S)}^2,
\]
because $\HermBasis_3(S)$ is one orthonormal degree-three Hermite direction. The trick we will use is to condition on the transverse variables, so define
\[
        A(s):=\E_T[h(s,T)].
\]
This is the average value of $h$ along the slice $S=s$. Since $h\in[-1,1]$,
\[
        |A(s)|\le1.
\]
Also, for every function $\varphi$ depending only on $S$, it is a straightforward conditioning to show that
\[
        \E[h(S,T)\varphi(S)]=\E[A(S)\varphi(S)].
\]
In particular,
\[
        \E[A(S)S]=\ip{h}{S}=\nu x
\]
and
\[
        \ip{h}{\HermBasis_3(S)}=\E[A(S)\HermBasis_3(S)].
\]

This is the same trick as in the two-dimensional note. The problem is now a variational problem in the one variable function $A$, even though we are working in $d$ dimensions. To bound the agreement part, it is sufficient to relax $A$, and answer the following question:

\begin{quote}
``Among all $A:\R\to[-1,1]$ with $\E[A(S)S]=\nu x$, how small can the cubic coefficient be?''
\end{quote}

The answer is the rearrangement lemma.

Define
\[
        W(S):=S^3-3S=\sqrt6\,\HermBasis_3(S).
\]
Also define
\[
        \Delta(x):=x+2(1+x)\log\frac{1+x}{2},
        \qquad
        \Delta_+(x):=\max\{\Delta(x),0\}.
\]

\begin{lemma}[One-dimensional rearrangement]
Let $S\sim N(0,1)$ and $x\in[0,1]$. Then
\[
        \sup_{\substack{|A|\le1\\ \E[A(S)S]=\nu x}}
        \E[A(S)(S^3-3S)]
        =-\nu\Delta(x).
\]
\end{lemma}

\begin{proof}
Fix $r\ge0$ by
\[
        e^{-r^2/2}=\frac{1+x}{2},
\]
and put
\[
        \mu=r^2-3.
\]
Let $W=S^3-3S$. Then
\[
        W-\mu S=S^3-3S-(r^2-3)S=S(S^2-r^2).
\]

For any feasible $A$,
\begin{align*}
        \E[AW]
        &=\E[A(W-\mu S)]+\mu\E[AS]\\
        &=\E[A(W-\mu S)]+\mu\nu x\\
        &\le \E|W-\mu S|+\mu\nu x,
\end{align*}
using $|A|\le1$.

Equality is attained by
\[
        A_*(S)=\sgn(W-\mu S)=\sgn(S)\sgn(S^2-r^2),
\]
provided $A_*$ has the right first moment. It does, because
\[
        A_*(S)S=|S|\sgn(S^2-r^2),
\]
so
\begin{align*}
        \E[A_*(S)S]
        &=2\E[|S|\one_{|S|>r}]-\E|S|\\
        &=2\nu e^{-r^2/2}-\nu\\
        &=\nu x.
\end{align*}
Thus $A_*$ is feasible and attains the dual upper bound.

It remains to compute its value. Since
\[
        \sgn(S)(S^3-3S)=|S|^3-3|S|,
\]
we get
\begin{align*}
        \E[A_*W]
        &=\E[(|S|^3-3|S|)\sgn(S^2-r^2)]\\
        &=2\E[(|S|^3-3|S|)\one_{|S|>r}]
          -\E[|S|^3-3|S|].
\end{align*}
The Gaussian tail moments are
\[
        \E[|S|\one_{|S|>r}]=\nu e^{-r^2/2},
        \qquad
        \E[|S|^3\one_{|S|>r}]=\nu(r^2+2)e^{-r^2/2}.
\]
Also
\[
        \E|S|=\nu,
        \qquad
        \E|S|^3=2\nu.
\]
Therefore
\begin{align*}
        \E[A_*W]
        &=2\nu(r^2-1)e^{-r^2/2}+\nu\\
        &=\nu\left(1+(1+x)(r^2-1)\right)\\
        &=\nu\left((1+x)r^2-x\right).
\end{align*}
Since $r^2=-2\log((1+x)/2)$,
\[
        \E[A_*W]
        =-\nu\left(x+2(1+x)\log\frac{1+x}{2}\right)
        =-\nu\Delta(x).
\]
This proves the lemma.
\end{proof}

% We delegate the proof to Appendix~\ref{app:rearrangement_lemma}, which uses standard variational techniques. Using this rearrangement lemma gives the desired bound on the agreement part.

\begin{corollary}[Diagonal bound]
For the general-dimensional ternary $h$ above,
\[
        \norm{P_3h}^2\ge \frac{\nu^2}{6}\Delta_+(x)^2.
\]
Consequently
\[
        2\norm{P_1h}^2-\norm{P_3h}^2
        \le
        \nu^2\left(2x^2-\frac16\Delta_+(x)^2\right).
\]
\end{corollary}

\subsection{The disagreement term}
\label{sec:disagreement}

Now we have to bound the disagreement term
\(
        \norm{P_3k}^2.
\)
For each fixed vector \(t\), restrict \(k\) to the one-dimensional line obtained by varying \(S\) while holding \(T=t\). We call this restriction the \textbf{fiber} of \(k\) over \(t\), and write
\[
    u_t(s):=k(s,t).
\]
Since \(k\in\{-1,0,1\}\), every fiber \(u_t\) is a one-dimensional ternary function. For $j=0,1,2,3$, define
\[
        K_j(t):=\E_S[k(S,t)\HermBasis_j(S)].
\]
Thus $K_j$ is the $j$th Hermite coefficient of the $S$-fiber, as a function of the transverse variable $t$. In two dimensions, the four degree-three basis elements were
\[
        \HermBasis_3(S),\qquad
        \HermBasis_2(S)\HermBasis_1(T),\qquad
        \HermBasis_1(S)\HermBasis_2(T),\qquad
        \HermBasis_3(T).
\]
In general dimensions, this becomes the orthogonal decomposition
\[
        \cH_3(S,T)
        =\bigoplus_{j=0}^3 \cH_j(S)\otimes \cH_{3-j}(T).
\]
This just says that a total degree-three Hermite can spend $j$ degrees in $S$ and $3-j$ degrees in $T$.

Therefore
\[
        P_3k
        =\sum_{j=0}^3 \HermBasis_j(S)\,\Pi^T_{3-j}K_j(T),
\]
where $\Pi^T_{3-j}$ is the orthogonal projection in the $T$ variable onto transverse Hermite degree $3-j$. Hence
\begin{equation}
        \norm{P_3k}^2
        =\sum_{j=0}^3 \norm{\Pi^T_{3-j}K_j}^2.
        \label{eq:p3k-decomposition}
\end{equation}
Since the $\Pi^T_{3-j}$ are orthogonal projections onto closed subspaces, they decrease norm, so we get the simpler upper bound
\[
        \norm{P_3k}^2
        \le \sum_{j=0}^3\norm{K_j}^2
        =\E_T\sum_{j=0}^3\left(\E_S[k(S,T)\HermBasis_j(S)]\right)^2.
\]
For a one dimensional ternary function $u$, define
\[
        V(u):=\sum_{j=0}^3\left(\E[u(Z)\HermBasis_j(Z)]\right)^2.
\]
Then the previous inequality becomes
\begin{equation}
        \boxed{\norm{P_3k}^2\le \E_T V(u_T).}
        \label{eq:p3k-fiber-bound}
\end{equation}

For a ternary function $u:\R\to\{-1,0,1\}$, define its weighted support budget
\[
        \beta(u):=\E[|Z|u(Z)^2].
\]
Since $u^2=\one_{u\neq0}$,
\[
        \beta(u)=\E[|Z|\one_{u(Z)\neq0}].
\]
So $\beta(u)$ is the $|Z|$-weighted Gaussian size of the support of the fiber.

\begin{theorem}[Fiber inequality]
For every ternary $u:\R\to\{-1,0,1\}$,
\begin{equation}
        V(u)\le 3\nu\beta(u)-\beta(u)^2.
        \label{eq:fiber-inequality}
\end{equation}
\end{theorem}


\paragraph{Proof sketch}. This is proved in Appendix~\ref{app:fiber-certificate}, which reduces ~\eqref{eq:fiber-inequality} to a finite computational check. Here, we briefly sketch out the proof. The inequality is genuinely finite-dimensional: $V(u)$ depends only on the four moments $\E[u\HermBasis_j]$, $0\le j\le3$. By the dual formula for the Euclidean norm,
\[
        \sqrt{V(u)}=\sup_{\norm{a}_2=1}\E[u(Z)p_a(Z)],
        \qquad
        p_a(Z)=\sum_{j=0}^3a_j\HermBasis_j(Z).
\]
For fixed $p_a$ and a fixed support price $\tau\ge0$, pointwise maximization over the three values of $u$ gives
\[
        \max_{v\in\{-1,0,1\}}\bigl(v\,p_a(z)-\tau|z|\,v^2\bigr)
        =\bigl(|p_a(z)|-\tau|z|\bigr)_+,
\]
with maximizer the threshold rule
\[
        u(z)=\sgn(p_a(z))\one_{\{|p_a(z)|\ge \tau |z|\}};
\]
threshold functions are thus the exact optimizers of the priced problem, not merely plausible candidates. The infinite-dimensional search over $u$ therefore collapses to the four-parameter family $(a,\tau)$, whose active sets are cut out by one-dimensional polynomial inequalities; the content of~\eqref{eq:fiber-inequality} is the inequality $3\nu\beta-m^2\ge\beta^2$ along this family, where $m=\int_E|p_a|\,d\gamma$ and $\beta=\int_E|z|\,d\gamma$ over the active set $E$. The constant fiber $u\equiv1$ (that is, $p_a=\HermBasis_0$, $\tau=0$) realizes $\rho:=(3\nu\beta-m^2)/\beta^2=3-\pi/2\approx1.4292$, and numerical scans over the family find no smaller value. This suggests that~\eqref{eq:fiber-inequality} --- which asserts only $\rho\ge1$ --- holds with about $0.43$ of true headroom, though we do not prove this sharper assertion.

Appendix~\ref{app:fiber-certificate} turns this into a proof: it reduces~\eqref{eq:fiber-inequality} to two certified regimes (a direction-free high-budget bound and a Fenchel-dual low-budget bound), each consisting of finitely many one-dimensional inequalities that are verified in outward-rounded interval arithmetic. 

\paragraph{Reproducibility.}
  The scripts, certificates, and Arb outputs for this 
  are available at \url{https://github.com/trishullab/grothendieck-bounds}. 

\vspace{\baselineskip}

Apply~\eqref{eq:fiber-inequality} to each fiber $u_T$. From~\eqref{eq:p3k-fiber-bound},
\[
        \norm{P_3k}^2
        \le
        \E_T\left[3\nu\beta(u_T)-\beta(u_T)^2\right].
\]
We now prove the following bound. Note that this step is dimension-free.

\begin{theorem}[Fiber Residual Bound \#1]
We have the fiber bound
\begin{equation}
        \norm{P_3k}^2
        \le
        \nu^2\left(3(1-x)-(1-x)^2\right).
        \label{eq:fiber-residual-one}
\end{equation}
\end{theorem}

\begin{proof}
We do this by relating the average budget to $x$. Since
\[
        \nu x=\E[hS]
\]
and
\[
        hS\le |S|h^2,
\]
we get
\begin{align*}
        \nu x
        &\le \E[|S|h^2]\\
        &=\E[|S|(1-k^2)]\\
        &=\nu-\E[|S|k^2].
\end{align*}
Therefore
\[
        \E[|S|k^2]\le \nu(1-x).
\]
But
\[
        \E[|S|k^2]
        =\E_T\E_S[|S|k(S,T)^2]
        =\E_T\beta(u_T).
\]
Let
\[
        b:=\E_T\beta(u_T).
\]
Then
\[
        b\le \nu(1-x).
\]

The function
\[
        \phi(t)=3\nu t-t^2
\]
is concave, so Jensen gives
\[
        \E_T\phi(\beta(u_T))\le \phi(\E_T\beta(u_T))=3\nu b-b^2.
\]
On $[0,\nu]$, $\phi$ is increasing. Hence, using $b\le \nu(1-x)$, we obtain~\eqref{eq:fiber-residual-one}.
\end{proof}

This is our first fiber residual bound. But this bound is not best for small $x$. We also use a second cruder bound.

\begin{theorem}[Fiber Residual Bound \#2]
    \begin{equation}
        \norm{P_3k}^2
        \le
        2\PhiG(\sqrt{-2\log x})-1.
        \label{eq:fiber-residual-two}
\end{equation}
\end{theorem}

\begin{proof}
Let
\[
        \rho(s):=\E_T[k(s,T)^2].
\]
Then $0\le\rho\le1$ and
\[
        \norm{k}^2_2=\E\rho(S).
\]
Also
\[
        \E[|S|\rho(S)]=\E[|S|k(S,T)^2]\le \nu(1-x).
\]
Since $P_3$ is an orthogonal projection,
\[
        \norm{P_3k}^2\le \norm{k}^2_2=\E\rho(S).
\]

Given a budget on $\E[|S|\rho(S)]$, the way to maximize $\E\rho(S)$ is to spend the budget where $|S|$ is smallest. Thus the extremal support is centered around zero.

Let
\[
        r_x:=\sqrt{-2\log x}
\]
for $x\in(0,1]$, with the convention $r_0=\infty$. Then
\[
        \E[|S|\one_{|S|\le r_x}]=\nu(1-e^{-r_x^2/2})=\nu(1-x).
\]
Therefore the bathtub principle gives
\[
        \norm{P_3k}^2
        \le
        \mathbb P(|S|\le r_x)
        =2\PhiG(\sqrt{-2\log x})-1,
\]
which proves~\eqref{eq:fiber-residual-two}.
\end{proof}

Combining the two fiber residual bounds, define
\begin{equation}
        R(x):=
        \min\left\{
        3(1-x)-(1-x)^2,
        \frac{1}{\nu^2}\left(2\PhiG(\sqrt{-2\log x})-1\right)
        \right\}.
        \label{eq:residual-function}
\end{equation}
Then
\[
        \norm{P_3k}^2\le \nu^2R(x).
\]

\subsection{The final inequality}
\label{sec:putting-together}

From the diagonal bound and the residual bound,
\[
        D_1(h,k)
        \le
        \nu^2\left(2x^2-\frac16\Delta_+(x)^2+R(x)\right).
\]
So it remains to prove the scalar inequality
\begin{equation}
        E_1(x):=2x^2-\frac16\Delta_+(x)^2+R(x)
        \le \frac{11}{6}
        \qquad(0\le x\le1).
        \label{eq:scalar-reduction}
\end{equation}

Let
\[
        r_-:=\frac{1-1/\sqrt3}{2},
        \qquad
        r_+:=\frac{1+1/\sqrt3}{2}.
\]
These are the roots of
\[
        x^2-x+\frac16=0.
\]

Now we do some simple casework based on the range in which $x$ lies. Everything is simple algebra.

\subsection*{Low range: $0\le x\le r_-$}

Use the crude bound $R(x)\le 1/\nu^2=\pi/2$ and drop the negative term:
\[
        E_1(x)\le 2x^2+\frac\pi2
        \le 2r_-^2+\frac\pi2.
\]
Now
\[
        2r_-^2+\frac\pi2
        =\frac23-\frac1{\sqrt3}+\frac\pi2
        <\frac{11}{6}.
\]

\subsection*{Middle range: $r_-\le x\le r_+$}

Use the fiber branch of $R$ and again drop the negative term:
\begin{align*}
        E_1(x)
        &\le 2x^2+3(1-x)-(1-x)^2\\
        &=x^2-x+2.
\end{align*}
Since $x\in[r_-,r_+]$,
\[
        x^2-x+\frac16\le0.
\]
Therefore
\[
        x^2-x+2\le \frac{11}{6}.
\]

\subsection*{High range: $r_+\le x\le1$}

Use the fiber branch again. We need the cubic correction.
It suffices to show
\begin{equation}
        \Delta(x)^2\ge 6x^2-6x+1.
        \label{eq:high-range-correction}
\end{equation}
Indeed, then
\begin{align*}
        E_1(x)
        &\le 2x^2+3(1-x)-(1-x)^2-\frac16(6x^2-6x+1)\\
        &=\frac{11}{6}.
\end{align*}

To prove~\eqref{eq:high-range-correction}, first note that
\[
        \Delta(x)-(3x-2)
        =2\left((1-x)+(1+x)\log\frac{1+x}{2}\right).
\]
Writing $t=(1+x)/2$, the bracket becomes
\[
        2\bigl((1-t)+t\log t\bigr)\ge0
\]
for $0<t\le1$. Hence
\[
        \Delta(x)\ge 3x-2.
\]
For $x\ge r_+$, both sides needed below are nonnegative, and
\[
        (3x-2)^2-(6x^2-6x+1)=3(1-x)^2\ge0.
\]
Thus
\[
        \Delta(x)^2\ge (3x-2)^2\ge 6x^2-6x+1.
\]
This proves the high range.

Combining the three ranges proves~\eqref{eq:scalar-reduction}.

We have thus shown
\[
        D_1(h,k)\le \frac{11}{6}\nu^2.
\]
As explained earlier, this implies the affine coefficient inequality
\[
        b_3\ge 2b_1-\frac{11}{6}.
\]
Since this is affine, it is preserved under mixtures.

\section{The Affine Coefficient Inequality Bounds the Inverse Majorant}
\label{sec:strip_majorant}

\begin{definition}[Admissible majorant value]\label{def:admissible-majorant}
Let $H$ be odd and analytic near the origin, with $H'(0)>0$. Let $H^{-1}$ denote the local inverse branch satisfying $H^{-1}(0)=0$, and write
\[
        H^{-1}(z)=\sum_{j\ge0}a_{2j+1}z^{2j+1}.
\]
A number $\gamma>0$ is \emph{admissible} for $H$ if this series converges absolutely at $\gamma$ and
\[
        \sum_{j\ge0}|a_{2j+1}|\gamma^{2j+1}\le1.
\]
\end{definition}

We now turn the affine coefficient inequality proved in Section~\ref{sec:affine-strip} into a universal barrier for the inverse-majorant parameter. The coefficients $b_1$ and $b_3$ depend linearly on the correlation function, so the strip is preserved under mixtures and coefficientwise limits. Section~\ref{sec:majorant_optimal} will apply this barrier to the limiting mixed functions supplied by the Naor--Regev construction.

\begin{proposition}[The majorant barrier]\label{prop:majorant-barrier}
Let
\[
        H(t)=b_1t+b_3t^3+b_5t^5+\cdots
\]
be odd and analytic near the origin. Suppose that
\[
        0<b_1\le1,
        \qquad
        b_3\ge2b_1-\frac{11}{6}.
\]
Then every admissible value for $H$ satisfies
\[
        \gamma\le\frac{11}{12}.
\]
\end{proposition}

\begin{proof}
Write
\[
        H^{-1}(z)=a_1z+a_3z^3+\cdots.
\]
Comparing the coefficients of $z$ and $z^3$ in $H(H^{-1}(z))=z$ gives
\[
        a_1=\frac1{b_1},
        \qquad
        a_3=-\frac{b_3}{b_1^4}.
\]

If $b_1\le11/12$, admissibility gives
\[
        1\ge |a_1|\gamma=\frac{\gamma}{b_1},
\]
so $\gamma\le11/12$.

Suppose now that $b_1>11/12$. The strip gives
\[
        b_3\ge2\left(b_1-\frac{11}{12}\right)>0.
\]
Let
\[
        M_H(s):=\sum_{j\ge0}|a_{2j+1}|s^{2j+1},
\]
allowing the value $+\infty$. Using only the first two terms,
\begin{align*}
        M_H\left(\frac{11}{12}\right)
        &\ge
        \frac{11}{12b_1}
        +\frac{2(b_1-\frac{11}{12})}{b_1^4}
          \left(\frac{11}{12}\right)^3\\
        &=1+
        \frac{(b_1-\frac{11}{12})
        \left(2(\frac{11}{12})^3-b_1^3\right)}{b_1^4}
        >1.
\end{align*}
The last inequality uses $b_1\le1$ and $2(11/12)^3>1$. If some $\gamma\ge11/12$ were admissible, monotonicity of $M_H$ would give
\[
        M_H\left(\frac{11}{12}\right)\le M_H(\gamma)\le1,
\]
a contradiction.
\end{proof}

\section{The Inverse Majorant Bound is Optimal}
\label{sec:majorant_optimal}

The key global input is obtained by reworking the construction in the proof of Naor and Regev's optimality theorem. We first record the notion of mixing needed to state it.

\begin{definition}[Mixed rounding scheme]\label{def:mixed-rounding}
Fix $d\in\mathbb N$. A $d$-dimensional \emph{mixed rounding scheme} is a probability distribution $\mu$ on pairs of measurable sign functions
\[
        f,g:\R^d\to\{\pm1\}.
\]
One pair $(f,g)$ is sampled from $\mu$ and then used for every vector in the rounding procedure. Its averaged correlation function is
\[
        H_\mu(t):=\E_{(f,g)\sim\mu}H_{f,g}(t).
\]
A \emph{limiting mixed correlation function} is a coefficientwise limit of averaged correlation functions of mixed rounding schemes, possibly in increasing dimensions.
\end{definition}

The admissible values from Definition~\ref{def:admissible-majorant} are the parameters appearing in the inverse-majorant bound~\eqref{eq:inverse-majorant-bound}. For an actual mixed scheme, the standard tensor preprocessing argument turns an admissible value $\gamma$ into the upper bound $K_G\le\pi/(2\gamma)$. The next lemma shows that, after allowing mixtures and increasing dimension, these bounds can approach $K_G$ itself.

\begin{lemma}[Optimality of the inverse-majorant bound]\label{lem:inverse-majorant-optimal}
There exist odd analytic limiting mixed correlation functions $H_k$ and admissible values $\gamma_k$ such that
\[
        \gamma_k\le\frac{\pi}{2K_G}
        \qquad\text{and}\qquad
        \gamma_k\longrightarrow\frac{\pi}{2K_G}.
\]
Equivalently, the resulting upper bounds satisfy
\[
        \frac{\pi}{2\gamma_k}\longrightarrow K_G.
\]
Thus, after allowing mixtures and increasing dimension, the inverse-majorant bound is asymptotically optimal.
\end{lemma}

\begin{proof}
We use the construction in the proof of Theorem~1.1 of \cite{naor2014krivine}. We do not use their theorem statement directly.

For each $k$, their equation~(4) gives a probability measure $\nu_k$ on pairs of sign assignments on $S^{k-1}$ whose averaged kernel is
\[
        \frac{1}{K_G}\ip{x}{y}.
\]
After radial extension, their equations~(5) and~(6) give an odd analytic normalized correlation function
\[
        H_k(t)=\frac{\pi}{2K_G}f_k(t),
\]
where $f_k$ is the function defined in their equation~(6). By partitioning the sphere into finitely many small Borel sets and extending the sampled signs as step functions, one obtains ordinary mixed rounding schemes whose averaged kernels converge uniformly to the kernel above. Hence $H_k$ is a limiting mixed correlation function in the sense of Definition~\ref{def:mixed-rounding}.

Naor and Regev write
\[
        f_k^{-1}(z)=\sum_{n\ge0}b_n(k)z^{2n+1}.
\]
Immediately before their Lemma~2.1, they choose a number $c_k\le1$ such that
\[
        \sum_{n\ge0}|b_n(k)|c_k^{2n+1}\le1,
\]
and Lemma~2.1 proves that $c_k\to1$. Define
\[
        \gamma_k:=\frac{\pi c_k}{2K_G}.
\]
Since
\[
        H_k^{-1}(z)=f_k^{-1}\left(\frac{2K_G}{\pi}z\right),
\]
the absolute inverse majorant of $H_k$ at $\gamma_k$ is exactly
\[
        \sum_{n\ge0}|b_n(k)|c_k^{2n+1}\le1.
\]
Thus $\gamma_k$ is admissible. Since $c_k\le1$ and $c_k\to1$,
\[
        \gamma_k\le\frac{\pi}{2K_G}
        \qquad\text{and}\qquad
        \gamma_k\longrightarrow\frac{\pi}{2K_G}.
\]
The equivalent statement about the resulting upper bounds follows by taking reciprocals.
\end{proof}

\begin{proof}[Proof of Theorem~\ref{thm:lower-bound}]
Let $H_k$ and $\gamma_k$ be supplied by Lemma~\ref{lem:inverse-majorant-optimal}. Sections~\ref{sec:preliminaries} and~\ref{sec:affine-strip} establish the affine coefficient inequality for every pair of odd sign functions in every dimension. It therefore holds for every mixed scheme, because $b_1$ and $b_3$ average linearly. Passing to the coefficientwise limit shows that
\[
        b_3(H_k)\ge2b_1(H_k)-\frac{11}{6}.
\]
This is the point at which preservation under mixing is used.

We also have $0<b_1(H_k)\le1$. Positivity follows because $\gamma_k$ is admissible. For the upper bound, every sign function $f$ satisfies
\[
        \norm{P_1f}
        =\sup_{\norm{u}=1}\left|\E[f(X)\ip{X}{u}]\right|
        \le\E|S|=\nu,
\]
so $|b_1(f,g)|\le(\pi/2)\nu^2=1$. This bound is also preserved under mixtures and coefficientwise limits.

Proposition~\ref{prop:majorant-barrier} now gives
\[
        \gamma_k\le\frac{11}{12}.
\]
Letting $k\to\infty$ and using Lemma~\ref{lem:inverse-majorant-optimal},
\[
        \frac{\pi}{2K_G}\le\frac{11}{12}.
\]
Rearranging proves the result.
\end{proof}
